\documentclass[11pt,a4paper]{article}
\usepackage[T1]{fontenc}
\usepackage{lmodern,amsmath,amssymb,amsthm,mathtools,booktabs,microtype}
\usepackage[margin=27mm]{geometry}
\usepackage[colorlinks=true,linkcolor=blue,urlcolor=blue,citecolor=blue]{hyperref}
\newtheorem{theorem}{Theorem}\newtheorem{lemma}[theorem]{Lemma}
\newtheorem{proposition}[theorem]{Proposition}
\newcommand{\Q}{\mathbb Q}\newcommand{\Z}{\mathbb Z}
\DeclareMathOperator{\Tr}{Tr}
\setlength{\emergencystretch}{2em}
\title{An explicit polynomial in twenty-nine variables\\
with effective elliptic sum heights}
\author{}\date{9 November 2026}
\begin{document}\maketitle
\begin{abstract}
Under the four arithmetic inputs of the specified elliptic construction,
we give an explicitly specified integer polynomial in twenty-nine total
variables whose rational input projection is exactly the integers.
Its total degree is less than $2^{89}$ and its coefficient height is
at most $2^{2^{99}}$. The construction uses direct square-abscissa
charts, a degree-six square predicate, and an effective height bound
for the sum of four elliptic abscissas. A single fixed coefficient
exponent, $137731$, is proved sufficient. There is no ineffective
constant in the specification. Every coefficient is determined by
integer arithmetic and finite determinant, resultant, and
product operations; independent programs verify the source identities,
all finite branches, and the degree and coefficient bounds.
\end{abstract}

\section{The defining polynomial and its arithmetic inputs}

All variable counts include the distinguished input. A polynomial
$F\in\Z[X,Y_1,\ldots,Y_m]$ defines $\Z$ over $\Q$ if
\[
 x\in\Z\quad\Longleftrightarrow\quad
 \exists\mathbf y\in\Q^m\quad F(x,\mathbf y)=0
 \qquad(x\in\Q).
\]
The coefficient height of an integer polynomial is the maximum
absolute value of its coefficients. The letter $h$ denotes the
absolute logarithmic Weil height, with $h(u/v)=\log\max(|u|,|v|)$
for relatively prime rational integers $u,v$ and $v>0$.

\begin{theorem}\label{thm:main}
Assume the four arithmetic assertions in Section~\ref{sec:inputs}.
The specification accompanying this paper determines a polynomial
\[
                   F\in\Z[X,Y_1,\ldots,Y_{28}]
\]
that defines $\Z$ over $\Q$ and satisfies
\[
             \deg F<2^{89},\qquad
             \max|\operatorname{coeff}(F)|\le2^{2^{99}}.
\]
Every coefficient choice, finite index range, and arithmetic operation
in the specification is explicit. In particular, the fixed odd exponent
$L=137731$ is a coefficient parameter, not a quantified variable.
\end{theorem}

The determinant specification is a compact exact description of an
ordinary polynomial. Matrix entries, finite-algebra generators, and
indices of finite products are not additional quantified coordinates.
The assumptions below remain explicit: the numerical and algebraic
checks do not replace their arithmetic proofs.

\subsection{The four inherited inputs}\label{sec:inputs}

Set
\[
 K=\Q(\sqrt2),\quad d=3-\sqrt2,\quad D=d^2=11-6\sqrt2,
 \quad E:y^2=x^3-Dx,
\]
\[
 P=\left(\frac{57}{8}-\frac72\sqrt2,
          -\frac{165}{8}+\frac{417}{32}\sqrt2\right).
\]
The construction uses the following assertions, with exactly the
scope established in the arithmetic dependency paper~\cite{Dependencies}.
\begin{enumerate}
\item $2E(K)=\Z P$, and $h(x([n]P))\ge n^2-6$ for every
      nonzero integer $n$.
\item For the five normalized contacts with
      $\infty,0,1,2401,-343/32$, the odd contact radical
      $M_5(q)$ outside primes at most $76176$ satisfies
      $h(q)\le H M_5(q)^{19}$ when $q$ is not a contact value,
      where
\[
 H=2^{110}210^{1800}\left(\prod_{p\le76176}p\right)^{19}+128.
\]
\item With $B=17b$, the pole system with preparation parameter
      zero is
\[
 p^2=B(1-e)(1+3e),\qquad
 Bq^2=(s^2-Be(B-b^2))(s^2+3Be(B-b^2)).
\]
      Its solvability excludes negative odd $v_\ell(b)$ for
      $\ell\notin\{2,3,17\}$. It is soluble whenever
      $0<b<17$, its reduced denominator is square, it is a
      square in $\Q_2,\Q_3,\Q_{17}$, and its valuations there
      are at least $3,1,2$, respectively. We retain the same
      restricted-converse and descent hypotheses supporting this
      assertion, and their simultaneous-realization construction.
      Lemma~\ref{lem:pole-chart} proves the nonzero-witness
      consequence of these inputs.
\item The simultaneous prime-pattern input permits the source's
      prescribed prime values, specified Gaussian generators, and reduction-order
      conditions, arbitrary fixed divisibility, and arbitrary
      fixed odd avoidance. It applies to the finite collection
      of comparison scales after the root points have been fixed.
\end{enumerate}
The prime-pattern input includes its stated application of Sun's
theorem. All four assertions have exactly the scope recorded in the
arithmetic dependency paper~\cite{Dependencies}. The effective
height lemma used here is proved in the companion paper~\cite{Effective};
it introduces no further arithmetic assumption.

\section{Effective control by the two components of a sum}

For four nonzero integer labels $n_1,\ldots,n_4$, put
\[
 Z=\sum_{i=1}^4x([n_i]P)=S_0+S_1\sqrt2,
 \qquad B=\max_i|n_i|.
\]
The effective elliptic sum theorem~\cite{Effective} proves
\begin{equation}\label{eq:aggregate}
                 B^2\le2h(Z)+80
                    \le2h(S_0)+2h(S_1)+83.
\end{equation}
Its proof uses formal-group estimates at finite places and positivity
at the real embeddings; it applies to all nonzero integer labels,
without a parity restriction. The underlying quantitative estimate is
\[
 h(Z)\ge B^2-24\log B-6-\log864.
\]
Indeed, $24\log B-B^2/2<27$ and $\log864<7$ imply the
first inequality in~\eqref{eq:aggregate}; the elementary height bound
$h(S_0+S_1\sqrt2)\le h(S_0)+h(S_1)+\tfrac32\log2$
implies the second. The companion proof treats cancellations at
all finite places and supplies the stated constants.

\begin{lemma}\label{lem:detector-height}
Let $H$ be the contact-height constant in Section~\ref{sec:inputs}.
Set $A_0=1$, $A_1=S_0^2$, and $A_2=(S_0+S_1)^2$.
Suppose the three comparison contact radicals at $A_r\lambda$
have their $k$th powers dividing a nonzero integer $N$.
Then
\begin{equation}\label{eq:commonheight}
 B^2<4(H+\log3)|N|^{19/k}+82,
 \qquad B^2-6<7H|N|^{19/k}.
\end{equation}
\end{lemma}
\begin{proof}
The abscissas of nonzero doubles are positive at both real
embeddings. Thus $S_0>\sqrt2|S_1|$, so $S_0>0$ and
$S_0+S_1>0$. All three scales are nonzero positive rational squares.
The fractional linear transformation $t\mapsto2+1/t$ identifies
the five comparison contacts with the five source contacts away
from the specified exceptional primes. The contact-height input gives
\[
 h(A_r\lambda)\le T:=(H+\log3)|N|^{19/k}\qquad(0\le r\le2).
\]
Taking quotients with the unscaled term gives $h(A_r)\le2T$,
and consequently
\[
                      h(S_0)\le T,\qquad h(S_0+S_1)\le T.
\]
Put $C=S_0+S_1$ and write
\[
                   Z=(1-\sqrt2)S_0+\sqrt2 C.
\]
The two coefficients are algebraic integers, so the nonarchimedean
local contribution is at most the sum of those of $S_0$ and $C$.
At the two real embeddings the additional factors are bounded by
$2\sqrt2-1$ and $2\sqrt2+1$, whose product is $7$. Thus
\[
                   h(Z)\le h(S_0)+h(C)+\tfrac12\log7
                              \le2T+\tfrac12\log7.
\]
Equation~\eqref{eq:aggregate} now gives
$B^2\le4T+80+\log7<4T+82$.
For the second bound use $|N|\ge1$, $H\ge28$, and
$4(H+\log3)+76<7H$.
\end{proof}

\section{The exact rational source}

We use the point charts, local tests, and denominator guards of
the established source~\cite{Baseline}. We describe their data
and the exact choices used here.

There are nine points over $K$, named
\[
 A,S,T_1,Z,H_0,U_A,U_S,U_T,U_Z.
\]
For each point retain two rational coordinates $h_0,h_1$ and set
\begin{equation}\label{eq:directpoint}
 h=h_0+h_1\sqrt2,\qquad x=h^2,\qquad
 j^2=h^4-D,\qquad y=hj.
\end{equation}
This chart describes precisely the nonzero points of $\Z P$.
Indeed, let $R_0=(1+\sqrt2,2\sqrt2)$, so $2R_0=P$.
The group input implies that every $Q\in E(K)$ is $mR_0+T$
with $T\in E[2](K)$. For the standard two-isogeny descent map,
$\delta(R_0)=1+\sqrt2$ has opposite signs at the real embeddings,
whereas the torsion classes are $1,-D,d,-d$, with equal signs.
A square abscissa therefore forces $m$ even. Of these torsion
classes only $1$ is square: $-D,-d$ are negative and
$N_{K/\Q}(d)=7$ is not a rational square. Thus $Q\in\Z P$.
Conversely, $\delta(2R)=1$ gives a square abscissa for every
nonzero multiple of $P$. Neither $h=0$ nor $j=0$ is possible
in~\eqref{eq:directpoint}. A direct proof of the descent map
and all these assertions is given in~\cite{SquareCharts}.

The finite root $j$ is represented by one monic quartic,
without a finite choice of charts. Write
\[
 h^4-D=A+B\sqrt2,
 \quad A=h_0^4+12h_0^2h_1^2+4h_1^4-11,
 \quad B=4h_0h_1(h_0^2+2h_1^2)+6.
\]
Set
\[
 c=2A+B,\qquad d_1=-A-B,\qquad j'=(c+d_1\sqrt2)j.
\]
Writing $(c+d_1\sqrt2)^2(A+B\sqrt2)=A'+B'\sqrt2$ gives
\[
 A'=6A^3-9AB^2-4B^3,\qquad
 B'=-4A^3+6AB^2+3B^3.
\]
For $j=u+v\sqrt2$, direct expansion gives
\[
 j'=2(u^3-2v^3)+(4v^3-u^3)\sqrt2.
\]
Its rational component is nonzero, since $2$ is not a rational
cube and $j\ne0$. The multiplier has norm $2A^2-B^2$,
which is nonzero for rational $(A,B)\ne(0,0)$ because $2$
is not a rational square. Thus the reconstruction uses only
these two elementary rational obstructions. The general quadratic
extension version is proved in~\cite{QuadraticCharts}.
We may therefore use the single chart
\begin{equation}\label{eq:pointquartic}
 Y^4-4A'Y^2+8(B')^2=0,\quad Y\ne0,\qquad
 j'=Y/2+(B'/Y)\sqrt2,\qquad
 j=\frac{j'}{c+d_1\sqrt2}.
\end{equation}
The point-chart factors $Y$, $A'^2-2B'^2$, $2A^2-B^2$,
and the quartic derivative are automatically nonzero at every
rational retained-base specialization and every algebraic root.
Indeed, $A+B\sqrt2=h^4-D$ cannot vanish: taking norms would
imply $N(h)^4=49$, which has no rational solution. The homogeneous cubic
$-4A^3+6AB^2+3B^3$ has no nontrivial rational zero, because
$4T^3-6T-3$ is Eisenstein at $3$. Finally,
\[
 A'^2-2B'^2=(2A^2-B^2)^2(A^2-2B^2)
\]
and the discriminant of the quartic in~\eqref{eq:pointquartic} is
\[
                 2^{15}B'^2(2A^2-B^2)^4(A^2-2B^2)^2\ne0.
\]
The exact square-test proof is given in~\cite{SquarePredicate}.
These four factors per point may therefore be omitted from the
finite-algebra localization product. Every other denominator of the
affine point operations remains guarded.

The four root labels $(n,a,t,b)$ are nonzero integers; no parity
condition is needed for soundness. The simultaneous completeness
construction retains its original odd labels. The scalar shift
is now computed directly as $S-P$, in the affine difference
chart. Its excluded equal-abscissa cases are $S=\pm P$, both
avoided by that construction.

Put $Q_0=P$, $(Q_1,Q_2,Q_3,Q_4)=(A,S,T_1,Z)$,
$U_0=H_0$, and $V_i=U_i-Q_i$. Write $z=-x/y$ and $f=1/x$.
The truncated formal logarithms are
\[
 \ell_r(Q)=z(Q)\sum_{j=0}^{r/4-1}c_j(Df(Q)^2)^j,
 \quad c_0=1,\quad c_1=-2/5,\quad
 c_j=c_{j-1}\frac{4j-5}{4j+1}.
\]
Use precisions $(12,40,40,40)$, ratio orders $(10,39,39,39)$,
and scalar order $20$. With
$\beta=216X^4+450459$, the sixteen tests are the ten
formal-group tests $f(U_i),\sigma f(V_i)$, the four differences
between the paired logarithmic ratios, the direct target test
for $\beta$, and the direct Lucas polynomial
\[
 v^2-u(u-2)vw+(u-2)^2w^2-1
\]
in the three corresponding ratios. Their assigned orders are
$1$ for each of the first ten tests and
$10,39,39,39,10,39$ for the remaining six. These tests and
their construction are unchanged.

\subsection{Three comparison scales}

We use precisely the three positive rational
square scales
\begin{equation}\label{eq:scales}
 A_0=1,\qquad
 A_1=S_0^2,\qquad A_2=(S_0+S_1)^2,
 \qquad S_0+S_1\sqrt2=\sum_{i=1}^4x(Q_i).
\end{equation}
They are nonzero: the root points are nonzero doubles and
their abscissas are positive in both real embeddings.
In fact $x(Q_i)>3-\sqrt2$ at each embedding, by the curve
equation and exclusion of two-torsion. Taking the sum and its
half-trace gives $S_0>12$, so $A_1\ne A_0$.
For the retained slope $\lambda$, each scale is inserted into
the same normalized comparison detector of~\cite{Baseline}:
\[
 \begin{gathered}
 t_r=A_r\lambda,\quad \delta_r=t_r-1,\quad
 F_r=\prod_{\theta\in\Theta}(t_r-\theta),\quad
 R_r=\frac{F_r}{A_r^5\delta_r^6},\\
 \Theta=\{0,-1/2,-1,1/2399,-32/407\}.
 \end{gathered}
\]
\[
 W_r=\prod_{j=1}^{16}\left(
 1+R_r\left(\frac{D_j}{\delta_r^{m_j+1}R_r^{m_j}}\right)^2\right),
 \qquad b_r=\frac{R_r}{\Tr(W_r^{-1})^2}.
\]
All displayed denominators and the inherited guards are retained,
except for the automatic point-chart factors proved above.
Each $b_r$ receives a pole certificate.

For each pole put $B=17b$. We use the single rational chart
\[
 e=\frac{\tau^2-B}{\tau^2+3B},\qquad
 p=\frac{-4B\tau}{\tau^2+3B},\qquad
 \tau^2+3B\ne0,
\]
with retained parameters $\tau,s$, and require the remaining quadratic
pole root $q$ to be nonzero. The following lemma proves completeness
of this restriction under the inherited arithmetic inputs.

\begin{lemma}\label{lem:pole-chart}
Suppose that $b$ has square reduced denominator,
\[
 0<b<17,\qquad b\in\Q_2^{\times2}\cap\Q_3^{\times2}
 \cap\Q_{17}^{\times2},\qquad
 v_2(b)\ge3,\quad v_3(b)\ge1,\quad v_{17}(b)\ge2.
\]
The inherited pole-realization proof supplies a certificate with
$pqs\ne0$ and $e\ne1$. Consequently the single chart above,
with $q\ne0$, suffices. The restricted predicate has the same
pole obstruction, and adds no arithmetic hypothesis.
\end{lemma}
\begin{proof}
Set $V=B(B-b^2)$. Its positivity and its local square conditions
are precisely those verified in the zero-preparation argument
of~\cite{Dependencies}. Let $\beta_0,\theta$ be the positive
squarefree representatives of $B,V$. Their supports are disjoint,
$\beta_0\equiv17\pmod{24}$, $\theta\equiv1\pmod{24}$, and
both mutual local-square conditions hold.

The inherited simultaneous-realization proof in~\cite{Glazer},
with source identifier
\begin{center}\small
\texttt{lem:original-seven-simultaneous-realization},
\end{center}
constructs points on $E_n:y^2=x(x-n)(x+3n)$ and $E_{n\theta}$ with nonzero
abscissas of squareclass $\beta_0$. Both are outside two-torsion.
Indeed, the construction includes a fresh prime of $n$ outside
$\beta_0\theta$, so $\beta_0$ differs in squareclass from $n$
and $n\theta$; its positive squareclass also differs from their
negative multiples by $-3$. These are the only nonzero two-torsion
abscissas. This non-torsion assertion is also explicit in the
source descent proof. Thus both ordinates are nonzero.

Writing $V=\theta c^2$ and rescaling as in that proof gives
\[
       (B,y_1)\in E_{e_0}(\Q),\qquad
       (Bs^2,y_2)\in E_{e_0V}(\Q),\qquad sy_1y_2\ne0.
\]
The substitutions $e=e_0/B$, $p=y_1/B$, $q=y_2/(B^2s)$
give exactly the two pole equations and $pqs\ne0$. The first
conic excludes $e=1$. Set $\tau=p/(e-1)$. Its equation gives
\[
       \tau^2+3B=\frac{4B}{1-e}\ne0,
\]
and then gives the displayed reconstruction of $e,p$.
Every restricted witness is an original pole witness, so its
obstruction is unchanged. The construction establishes completeness
for every input satisfying the displayed conditions. No equivalence
of the two pole predicates for arbitrary rational $b$ is asserted.
\end{proof}

\subsection{A linear scalar normalization}

The fixed scalar normalization is
\begin{equation}\label{eq:normalizer}
 \begin{gathered}
 \ell=76249,\quad
 r=\beta(\beta-1)(\beta-2401)(32\beta+343)(2\beta-1),\\
 R_L=\frac{102^2\ell^L r}{(1+C_Lr)^2},\quad
 C_L=17(\ell^{(L+1)/2}+1).
 \end{gathered}
\end{equation}
Here $L=137731$ is fixed, and thus $(L+1)/2=68866$. The scalar norm tree and its six rational
component tests are unchanged: they concern
$f(H_0),f(S-P),z(S-P)/z(H_0)$ with orders
$1,1,1,1,20,20$. Its output $W$ supplies the scalar pole
input $R_LW^2$.

The finite residue certificate verifies that the five contact
values are not attained by $216x^4+450459$ modulo $\ell$,
and that $1+17r$ never vanishes there. Consequently, for every
rational input,
\begin{equation}\label{eq:valuation}
 v_\ell(R_L)=
 \begin{cases}
 L,&v_\ell(X)\ge0,\\
 L-20v_\ell(X),&v_\ell(X)<0.
 \end{cases}
\end{equation}
These valuations are positive and odd. At a consumed contact
prime different from $\ell$, the denominator is a unit and the
contact valuation is unchanged. For integral $X$, the reduced
denominator of $R_L$ is square, it has the required local
squareclasses and depths at $2,3,17$, and
\[
 0<R_L<\frac{36}{\ell r}<1.
\]
The scalar norm argument therefore gives
\begin{equation}\label{eq:depth}
 v_\ell(a-1)\ge19L+1.
\end{equation}
For completeness, if $v_\ell(X)<0$, the leading coefficient of
$r$ is an $\ell$-unit, so $v_\ell(r)=20v_\ell(X)<0$ and
$v_\ell(1+C_Lr)=v_\ell(r)$, proving~\eqref{eq:valuation}.
If $v_\ell(X)\ge0$, finite-field avoidance makes both $r$
and $1+C_Lr$ units. At a consumed finite contact prime
$p\ne\ell$, one has $v_p(r)>0$, so $1+C_Lr$ is a unit.
The real bound follows from
$(1+C_Lr)^2>C_L^2r^2>17^2\ell^{L+1}r^2$.
For integer input, the denominator is an integer square;
$\ell$ and $102^2$ preserve the three required local squareclasses,
and the normalizer is a unit at $2,3,17$. The inherited scalar
norm argument applied at the forced contact $\ell$ gives
\eqref{eq:depth}. The pole preparation and simultaneous realization
are the same as in~\cite{Baseline}. The supplied
checker independently verifies every residue modulo $76249$
and the normalization identity as an exact polynomial identity.

\section{Soundness with the fixed exponent}

Let a rational source solution be given. The point charts yield
nonzero integer labels $(n,a,t,b)$ for the four root points.
The affine guard for $S-P$ excludes $a=1$. Put
\[
 c=|a-1|,\qquad B=\max(|n|,|a|,|t|,|b|),
 \qquad\Delta=t^2-a(a-2)tb+(a-2)^2b^2-1.
\]
The fixed-prime depth gives
\[
 c\ge\ell^{2616890},\qquad
 16\cdot2616890=41870240>41869080=20(2093448+6).
\]
Since $\ell>2^{16}$ and $H<2^{2093448}$, it follows that
\begin{equation}\label{eq:cutoff}
                         c>(64H)^{20}.
\end{equation}
In particular $B\ge c-1>4$, and $|\Delta|\le5B^4$.
The unchanged formal-logarithm and scalar-norm arguments
of~\cite{Baseline} give
\begin{align}
 h(\beta)&\le Hc^{19/20},\label{eq:betaheight}\\
 \Delta\ne0&\Longrightarrow
 B^2-6<7H(5B^4)^{19/39},\label{eq:delta}\\
 \beta\ne n&\Longrightarrow
 B^2-6<7H(2e^{h(\beta)}B)^{19/10}.\label{eq:beta}
\end{align}
For the last two implications, the local tests give the required
integer divisibilities and Lemma~\ref{lem:detector-height}
supplies their height consequence. The scalar inequality is
unchanged because the contact valuations consumed outside $\ell$
are unchanged by the linear normalizer.

Equation~\eqref{eq:cutoff} implies $h(\beta)<c/64$.
If $\Delta\ne0$, use $5^{19/39}<\sqrt5<16/7$ and
$B^2>12$ in~\eqref{eq:delta} to obtain
\[
                        B^{2/39}<32H.
\]
On the other hand,
\[
 B\ge c-1>c/2>(64H)^{20}/2>(32H)^{39/2},
\]
a contradiction. Thus $\Delta=0$.
The elementary Lucas growth lemma of~\cite{Baseline} applies
to these nonzero integer labels and gives $\log B>c-4$.
Its hypotheses require $|a|>7$, which follows from~\eqref{eq:cutoff}.

If $\beta\ne n$, equation~\eqref{eq:beta} gives
\[
 \log B<10\log(64H)+19h(\beta)
            <29c/64<c-4.
\]
Here $\log(64H)<c/64$ follows from $H\ge1$
and~\eqref{eq:cutoff}. This contradicts the Lucas growth bound.
Consequently $\beta=n\in\Z$.
If $X=u/v$ is in lowest terms with $v>0$, integrality of
$216X^4+450459$ implies $v^4\mid216=2^3 3^3$.
Hence $v=1$ and $X\in\Z$. This proves soundness.

\section{Completeness}

Fix an integer input $X$ and the stated fixed exponent $L$. All numerator
and denominator primes of~\eqref{eq:normalizer} form a finite
set. The source realization first chooses the odd root labels
with the scalar label $a-1$ sufficiently divisible and deep at
these primes, and with the corresponding points in the required
real neighborhoods. The explicit Lucas construction then gives
odd nonzero $t,b$ with $\Delta=0$ and arbitrarily large
growth. The input label is chosen to be $n=\beta$.

At this stage the three scales in~\eqref{eq:scales} are fixed
positive rational squares. Remove repetitions before applying the
simultaneous prime-pattern theorem. At least two distinct scales
remain, since $S_0>12$; thus its hypothesis requiring at least two
scales is satisfied. Equal scales use the same detector and pole
witnesses in the source. The theorem chooses the slope
and contact primes with the prescribed conjugate reduction
orders while avoiding the previously fixed odd modulus. The
subsequent Chinese remainder choice of the auxiliary index $h$
and the points $U_i=[hn_i]P$ supplies the sixteen local tests
to their assigned precisions and preserves the fixed scalar
requirements and real neighborhoods. This is precisely the
simultaneous construction of~\cite{Dependencies,Baseline}; its
prime-pattern input accepts any fixed finite collection of the
positive square scales and any fixed finite avoidance modulus.

Each of the three comparison inputs is thereby prepared for the
zero-preparation pole theorem. The scalar input is prepared as
well: $R_L$ has square reduced denominator and the required
local squareclasses, the six normalized scalar tests can be made
arbitrarily small at the fixed finite places, and the scalar
norm output can be made arbitrarily close to $1$ in the real
absolute value. Thus $0<R_LW^2<17$, with the requisite
denominator and local conditions. The four pole certificates
exist by the inherited completeness input. The direct point chart~\eqref{eq:pointquartic} covers all nine
chosen points; its multiplier needs no further choice.
Lemma~\ref{lem:pole-chart} applies independently to each pole input.
Its parameters belong only to that pole, so the certificates can
be chosen simultaneously without any common-parameter assumption.
All pole roots are nonzero and the single conic chart applies.
All remaining inherited guards are maintained by the same realization.

This proves completeness for the fixed specification. Together with
the preceding section it proves that the rational input projection
of the source is exactly $\Z$.

\section{Finite-algebra elimination and the exact bounds}

The retained rational coordinates are
\begin{center}
\begin{tabular}{lr}
\toprule
Coordinates & Number\\
\midrule
Input & 1\\
Nine square-root-abscissa pairs & 18\\
Comparison slope & 1\\
Three comparison-pole parameter pairs & 6\\
Scalar-pole parameter pair & 2\\
\midrule
Total retained before elimination & 28\\
\bottomrule
\end{tabular}
\end{center}
There are nine monic quartic point roots and four monic
quadratic pole roots. The first conics use one chart, and every
quadratic pole root is already guarded nonzero in the source.
Thus zero-root strata are inconsistent and are discarded.
The finite-algebra rank is $4^9 2^4=2^{22}$.
The point quartics are separable at every rational retained base.
The theorem of~\cite{Active} replaces all thirteen finite roots
by one rational coordinate while retaining every nonautomatic guard.
There are precisely 29 total polynomial variables.

Write $B_0$ for the algebra of the nine point roots and $A_0$
for the algebra after the quadratic pole roots have been adjoined.
These algebra symbols are unrelated to the earlier point-chart
coefficients. Put $b_0=4^9=2^{18}$ and $n_0=\dim A_0$.
If the pole roots are $q_1,\ldots,q_s$, their relations are
$q_j^2=c_j$ with $c_j\in B_0$. Exact coefficient-dependency
inspection shows that the source guard is
$G=gq_1\cdots q_s$ with $g\in B_0$. There are no closed
residual equations. Define $T=g^2\prod_jc_j$ and
\[
 \det_{B_0}(I-vM_T)=\sum_{j=0}^{b_0}a_jv^j,\qquad a_0=1.
\]
Let $\theta_1,\ldots,\theta_m$ be the point and pole generators
in order, and let $\delta_i$ be the discriminant of the $i$th
point quartic. They are scalar polynomials, nonzero at every
rational base point by the point-chart proof above. Put
\[
 R=n_0\left(1+\sum_{i=1}^m\sum_{u,v=1}^{n_0}
          (M_{\theta_i}^{A_0})_{uv}^2
        +\sum_{u,v=1}^{b_0}(M_g^{B_0})_{uv}^2\right).
\]
For each $1\le r\le b_0$, set
\begin{align*}
 c&=a_r,& h_r(Z)&=\sum_{j=0}^r a_jZ^{r-j},& E_r&=c-h_r(T),\\
 D_r&=c^2\prod_{i=1}^9\delta_i^2,&
 S_r&=r(s+1)+108,\\
 K_r&=(2R)^{S_r}(1+c^2)\prod_{i=1}^9(1+\delta_i^2),&
 C_r&=D_r+4RK_r,\\
 \beta_r&=D_rE_r\left(C_r^m+
                  \sum_{i=1}^mD_r^iC_r^{m-i}\theta_i\right),&
 H_r(U)&=\det_{A_0}(I-UM_{\beta_r}).
\end{align*}
The required polynomial is the following finite integer product:
\[
 F=H_{b_0}(U)\prod_{r=1}^{b_0-1}
             \left(H_r(U)^2+\sum_{j=r+1}^{b_0}a_j^2\right).
\]
Every operation has an explicit finite range. There is exactly
one extra rational coordinate, namely $U$.

The separating-element theorem of~\cite{Active} proves that
this polynomial has precisely the guarded source projection.
Here is the reason for its treatment of every specialization.
The actual degree of $\det_{B_0}(I-vM_T)$ counts the point-algebra
characters with $T\ne0$. At that degree, $E_r/c$ is the
idempotent of those factors. Moreover
\[
 |c|=\prod_{\chi:\,T(\chi)\ne0}
                     |g(\chi)|^2\prod_j|c_j(\chi)|.
\]
Since all generator and $g$-values have absolute value at most
$R$, isolating a factor bounds each required quadratic root
separation from below. The scalar discriminants similarly bound
the point-root separations. Thus $D_r/K_r$ is a valid positive
separation bound, and the powers of $C_r/D_r$ distinguish the
characters where the guard is nonzero. Multiplication by $E_r$
annihilates the other factors, including their nilpotents.
A rational root of $H_r$ therefore gives a one-dimensional
rational simultaneous eigenspace of all multiplication matrices,
which reconstructs a guarded rational character. Conversely
any such character gives a rational root. If $c=0$, then
$H_r=1$; the squared tail coefficients exclude every other
incorrect degree case. This also covers complex residue factors
and an empty guarded part. The full proof and coefficient
estimates are in~\cite{Active}.

For clarity, if the certified weight and logarithmic operator
norm of $\beta_r$ are at most $W_1r+W_0$ and $L_1r+L_0$,
respectively, the principal-minor estimates give
\begin{align*}
 \deg F&\le n_0\bigl(W_1b_0^2+(W_0+1)(2b_0-1)\bigr),\\
 \log_2\|F\|_1&\le n_0\bigl(L_1b_0^2+
                   (L_0+\lceil\log_2n_0\rceil+1)(2b_0-1)\bigr)\\
 &\hspace{12mm}+(b_0-1)\lceil\log_2(b_0+1)\rceil.
\end{align*}
The proof sums every degree case and uses
$2\sum_{r=1}^{b_0-1}r+b_0=b_0^2$.

Here $m=13$, $s=4$, and $n_0=2^{22}$. The point-generator
weights are all 6; the remaining weights are
$1568138,1935231,1935231,116318$. The complete guard and
its point factor have weights $6789625$ and $1234707$.
The exact linear majorants are
\[
 W_1=1367538680,\quad W_0=21229494749,
 \qquad L_1=1937703183919,\quad L_0=33390373627586.
\]
They give
\[
 \deg F\le394212872202031272474705920<2^{89},
\]
\[
 \log_2\|F\|_1\le558578343639412157650224021485<2^{99}.
\]
Since coefficient height is at most $\|F\|_1$, these prove
the stated bounds. They are upper bounds on the specified
polynomial, not claims about its exact expanded degree or height.
All finite sums and integer inequalities are independently checked.

\section{Reproduction and computational scope}

The directory \texttt{effective29\_nontorsion} contains the exact
compressed source and all programs for its construction and independent
source verification. From that directory,
\begin{verbatim}
python3 reproduce.py
\end{verbatim}
reconstructs the source at the proved sufficient exponent $137731$.
Its source has $7944$ coefficient nodes, and its integer coefficient
graph has $3243$ nodes. Checks cover nine quartic point charts, four
rational-square clearings, all $114$ atomic source guard factors,
and the $36$ automatic factors omitted from the guard. The scalar
normalizer is checked at all $76249$ residue classes, and its cutoff
inequalities are checked by exact integer comparisons.

\begin{samepage}
The directory \texttt{active\_elimination} contains the compiler and
independent checker:
\begin{verbatim}
build_active_specification.py
check_active_specification.py
check_active_degenerations.py
\end{verbatim}
\end{samepage}
The final specification is in \texttt{active29}. The compiler uses the verified source
coefficient graph and writes precisely the polynomial defined above.
The checker independently verifies the source dependencies and guard,
certifies the companion-column inequalities, and sums every degree case
literally. A separate checker uses determinant expansions, Newton traces,
and exact polynomial gcds on small finite algebras with complex residue
factors, inactive nilpotents, empty guarded parts, and degree drops.
All programs require only the Python standard library.

These computations verify the specified algebra and numerical bounds.
The nonvanishing arguments, effective height theorem, and four
arithmetic inputs have the proofs and scope stated above. Exact finite
coefficient reconstruction does not prove those global arithmetic
premises.

\begin{thebibliography}{99}
\bibitem{Active}
\emph{Separating the Nonzero Part of a Finite Algebra},
companion paper, 9 November 2026,
\nolinkurl{active_finite_algebra_20261109_v15.tex}.

\bibitem{Effective}
\emph{Effective heights of sums of elliptic abscissas},
companion paper, 8 November 2026,
\nolinkurl{effective_elliptic_sums_20261108_v10.tex}.
\bibitem{Baseline}
\emph{An explicit polynomial in 29 variables over the rationals},
the elliptic source, Lucas growth lemma, and guarded finite-algebra
construction, version of 8 November 2026,
\nolinkurl{few_variables_20261108_v8.tex}.
\url{https://github.com/JoseBrox/diophantine-z-q_improvements}.
\bibitem{Dependencies}
\emph{Arithmetic dependencies of the construction with 29 variables},
version of 8 November 2026,
\nolinkurl{arithmetic_dependencies_20261108_v8.tex}.
\url{https://github.com/JoseBrox/diophantine-z-q_improvements}.
\bibitem{SquareCharts}
\emph{Square abscissas and direct charts for the elliptic subgroup},
companion paper, 8 November 2026,
\nolinkurl{square_abscissa_charts_20261108_v9.tex}.
\bibitem{QuadraticCharts}
\emph{A single quartic chart for square roots in a quadratic extension},
companion paper, 8 November 2026,
\nolinkurl{quadratic_square_chart_20261108_v9.tex}.
\bibitem{SquarePredicate}
\emph{Optimal polynomial square tests in quadratic number fields},
companion paper, 8 November 2026,
\nolinkurl{quadratic_square_predicates_20261108_v10.tex}.
\bibitem{Glazer}
E.~Glazer, \emph{Diophantine definition of the integers over the
rationals}, Version~9, \nolinkurl{manuscripts/quartic_v9.tex},
commit \nolinkurl{14e8f601fe0d9173d4a5f084aedc3e07066a7688}.
\url{https://github.com/elliotglazer/diophantine-z-q}.
The inherited restricted-converse correction is identified
in~\cite{Dependencies}.
\end{thebibliography}
\end{document}
